\section{问题四：独立任务分区与资源配置}
\noindent{\heiti 版本说明：}本章已按最新Q3的22运输架次、3中继架次、6926.1512~s方案重算；运输与中继任务、通信保障关系均保持不变。

\subsection{任务继承与分区约束}

问题四严格继承问题三的货箱组批、访问顺序、起止时刻及通信保障关系，不再重排运输或中继任务。以15个服务区为顶点；若两个服务区出现在同一运输架次中，则将其强制合并，且合并关系具有传递性。新方案形成两个多站单元：\{S002,\allowbreak S009\}与\{S003,\allowbreak S007,\allowbreak S011,\allowbreak S015\}；其余9站各自独立，合计11个不可拆单元。

设服务区分组为$G_1,\ldots,G_K$，其中$K\in\{2,3\}$。各组非空、两两不交且并集为全部服务区。某组继承所有涉及该组服务区的运输架次；若某中继架次曾保障该组任一运输架次，则本模型为该组保留该中继的完整任务窗口。跨组共享同一中继时，在各组独立复制所需保障，不跨组调配资源。原题强制同组条款针对运输架次，并未要求同一中继保障的所有服务区合并；上述完整窗口复制是本章明确采用的保守核算假设，不等同于允许共享中继实体或截短原中继安排。

\subsection{峰值资源与分区目标}

对组$g$和机型$a$，运输无人机需求为任务区间的最大重叠数
\[
N^{\rm U}_{ga}=\max_t\sum_{r\in\mathcal R_{ga}}
\mathbf 1\{s_r\le t<e_r\}.
\]
共享电池从架次开始至返航后充满均视为占用，将$e_r$替换为充电完成时刻即可得到$N^{\rm B}_{ga}$。中继无人机占用至返航后300~s周转结束，能源组件占用至充满；两类中继资源同样按区间最大重叠数核算。时刻统一保留到0.001~s，避免输出舍入造成伪重叠。组内同型资源可以重新编号但不改变任务时刻，区间着色可实现峰值配置；因此资源需求不是Q3表中不同编号的机械计数。

组间工作量定义为该组全部运输与中继架次从开始至返航的时长之和，以变异系数
\[
CV=\frac{\sqrt{K^{-1}\sum_g(W_g-\bar W)^2}}{\bar W}
\]
衡量均衡性。沿用前轮建模选择，以$CV\le0.10$作为均衡目标，再依次最小化相对现有库存的资源总缺口、$CV$、货箱数极差、重复配置规模及空间离散度。10\%阈值不是题面硬约束；各类资源数量直接相加也不表示采购成本加总。对11个单元的$S(11,2)=1023$个两组分区和$S(11,3)=28501$个三组分区穷举，两者均有满足均衡目标的解。下述最优性仅针对冻结的Q3、完整中继窗口复制及给定目标次序成立。

\subsection{分区与资源配置结果}

\begin{table}[htbp]
\centering
\caption{两组任务分区及独立资源需求}
\small
\begin{tabular}{c p{5.2cm} *{8}{>{\centering\arraybackslash}p{0.8cm}}}
\toprule
组别 & 服务区 & \shortstack{A\\机} & \shortstack{B\\机} & \shortstack{C\\机} & \shortstack{A\\电池} & \shortstack{B\\电池} & \shortstack{C\\电池} & \shortstack{中继\\机} & 组件\\
\midrule
G1 & \shortstack[l]{S001,S003,S005,S006,\\S007,S011,S012,S015} & 3&1&2&4&1&3&2&2\\
G2 & \shortstack[l]{S002,S004,S008,S009,\\S010,S013,S014} & 2&2&1&3&3&1&2&2\\
\bottomrule
\end{tabular}
\end{table}

两组各承担11个运输架次；G1含49箱、2个中继任务副本，G2含31箱、3个中继任务副本。两组综合工作量分别为29844.080~s和31492.948~s，$CV=0.026882$。合计需A/B/C型运输机$5/3/3$架、对应电池$7/4/4$组、中继机4架和能源组件4组。

\begin{table}[htbp]
\centering
\caption{三组任务分区及独立资源需求}
\small
\begin{tabular}{c p{5.2cm} *{8}{>{\centering\arraybackslash}p{0.8cm}}}
\toprule
组别 & 服务区 & \shortstack{A\\机} & \shortstack{B\\机} & \shortstack{C\\机} & \shortstack{A\\电池} & \shortstack{B\\电池} & \shortstack{C\\电池} & \shortstack{中继\\机} & 组件\\
\midrule
G1 & \shortstack[l]{S001,S003,S005,\\S007,S011,S015} & 3&1&2&4&1&2&1&1\\
G2 & S002,S004,S009,S013 & 2&1&1&3&1&1&2&2\\
G3 & \shortstack[l]{S006,S008,S010,\\S012,S014} & 0&2&1&0&3&1&2&2\\
\bottomrule
\end{tabular}
\end{table}

三组分别承担9、7、6个运输架次，货箱数为40、20、20，中继任务副本数为1、3、2，综合工作量为22791.101、24197.998、19648.754~s，$CV=0.085615$。合计需A/B/C型运输机$5/4/4$架、对应电池$7/5/4$组、中继机5架和能源组件5组。

\subsection{两种方案比较与缺口原因}

\begin{table}[htbp]
\centering
\caption{两种分区方案的资源与均衡性比较}
\small
\begin{tabular}{crrrrrr}
\toprule
$K$ & 配置总数 & 重复配置数 & 库存缺口数 & 库存冗余数 & 工作量$CV$ & 货箱极差\\
\midrule
2 & 34 & 8 & 6 & 2 & 0.0269 & 18\\
3 & 39 & 13 & 10 & 1 & 0.0856 & 20\\
\bottomrule
\end{tabular}
\end{table}

现有库存为A/B/C型运输机$4/2/2$架、电池$6/4/4$组、中继机2架和能源组件6组。两组方案缺A/B/C型运输机各1架、A型电池1组和中继机2架，合计6个资源单位；三组方案缺A型机1架、B/C型机各2架、A/B型电池各1组和中继机3架，合计10个资源单位。两组与三组方案分别富余2组和1组中继能源组件，异类库存富余不能抵消其他类型的缺口。

不分区时，同一Q3任务的峰值需求为运输机$4/2/2$、电池$6/4/4$、中继机2架和组件2组，合计26个资源单位；虽然Q3方案使用了3个组件编号，按不变时间窗核算的峰值只需2组。分为两组和三组后分别重复配置8个和13个资源单位。

缺口主要来自跨组复用的丧失及中继保障的独立复制，而非运输任务增加。新Q3整体更省能、更快，并不保证分区后更省资源，因为多站批次和资源复用关系也发生了改变。本轮两组方案不仅配置规模和库存缺口较小，工作量$CV$及货箱极差也小于三组方案；三组提供更细的地域责任划分，但没有表现出更好的工作量均衡。因而在上述目标与核算口径下推荐两组方案；其执行仍以补齐所列缺口为前提。

\subsection{独立检验}

独立校验程序从Q3结果重新建立运输和中继占用区间，核对服务区互斥完备划分、多站批次不可拆、任务继承、资源峰值、逐项库存缺口及冗余、运输和中继能耗、综合工作量。两组与三组方案全部通过复核；另以三个回归测试检查同刻周转、充电重叠及同组约束的传递性。Q3输入快照及SHA-256清单随结果保存，确保上下游版本可追溯。官方结果表填入分区及资源数量，完整比较另存明细文件。
